\documentclass[aspectratio=169,11pt]{beamer}
\usetheme{Madrid}
\usecolortheme{dolphin}
\setbeamertemplate{navigation symbols}{}
\setbeamertemplate{caption}[numbered]

\usepackage[utf8]{inputenc}
\usepackage[T1]{fontenc}
\usepackage[spanish,es-noshorthands]{babel}
\usepackage{amsmath,amssymb,amsthm,mathtools}
\usepackage{tikz}
\usepackage{pgfplots}
\pgfplotsset{compat=1.18}
\usepackage{booktabs}
\usepackage{array}

\title[Prob. y Estadística]{Clase 18: Teorema central del límite}
\subtitle{Unidad 4: Función generadora de momentos}
\institute[UNS]{Universidad Nacional del Sur \\ Departamento de Matemática}
\author{Probabilidad y Estadística (5765)}
\date{}

\begin{document}

\begin{frame}
\titlepage
\end{frame}

\begin{frame}{Contenidos}
\tableofcontents
\end{frame}

\section{Motivación}

\begin{frame}{De la ley de los grandes números al TCL}
\begin{block}{Lo que ya sabemos (Clase 17)}
Si $X_1,\dots,X_n$ son i.i.d.\ con $E(X_i)=\mu$, $\text{Var}(X_i)=\sigma^2$, entonces $\overline{X}_n \xrightarrow{P} \mu$: el promedio muestral se acerca a $\mu$, y su varianza $\sigma^2/n \to 0$.
\end{block}

\begin{alertblock}{Pregunta natural}
Sabemos que $\overline{X}_n - \mu$ tiende a $0$. Pero, ¿con qué \textbf{forma} (distribución) se acerca a $0$? Si ``agrandamos la lupa'' apropiadamente (reescalando por $\sqrt n$), ¿qué distribución límite aparece?
\end{alertblock}

Este es el contenido del \textbf{teorema central del límite} (TCL), posiblemente el resultado más importante de toda la teoría de probabilidad clásica, por su enorme generalidad y sus aplicaciones prácticas.
\end{frame}

\section{Enunciado del TCL}

\begin{frame}{Estandarización}
\begin{block}{La variable estandarizada}
Sean $X_1,\dots,X_n$ i.i.d.\ con $E(X_i)=\mu$, $\text{Var}(X_i)=\sigma^2 < \infty$, y sea $S_n = X_1+\cdots+X_n$. Sabemos:
\[
E(S_n) = n\mu, \qquad \text{Var}(S_n) = n\sigma^2.
\]
Definimos la versión estandarizada de $S_n$ (o, equivalentemente, de $\overline{X}_n$):
\[
Z_n = \frac{S_n - n\mu}{\sigma\sqrt{n}} = \frac{\overline{X}_n - \mu}{\sigma/\sqrt{n}}.
\]
Por construcción (Clase 15), $E(Z_n) = 0$ y $\text{Var}(Z_n) = 1$ para todo $n$.
\end{block}
\end{frame}

\begin{frame}{Enunciado}
\begin{theorem}[Teorema central del límite]
Sean $X_1,\dots,X_n$ variables aleatorias i.i.d.\ con $E(X_i)=\mu$ y $\text{Var}(X_i)=\sigma^2 \in (0,\infty)$. Entonces, cuando $n\to\infty$, la distribución de
\[
Z_n = \frac{\overline{X}_n - \mu}{\sigma/\sqrt{n}} = \frac{S_n - n\mu}{\sqrt{n}\,\sigma}
\]
converge a la distribución normal estándar $N(0,1)$. Es decir, para todo $z \in \mathbb{R}$:
\[
\lim_{n\to\infty} P(Z_n \leq z) = \Phi(z) = P(Z\leq z), \quad Z\sim N(0,1).
\]
\end{theorem}

\begin{alertblock}{Lo notable}
El resultado vale sin importar cuál sea la distribución \textbf{original} de los $X_i$ (discreta, continua, simétrica, asimétrica...), siempre que tenga media y varianza finitas. Por eso se llama ``central'': es un resultado unificador.
\end{alertblock}
\end{frame}

\begin{frame}{Idea de la demostración usando la FGM}
\begin{block}{Estrategia}
Usamos la unicidad de la función generatriz de momentos (Clase 16): si mostramos que $M_{Z_n}(t) \to M_Z(t) = e^{t^2/2}$ (la FGM de $N(0,1)$) para todo $t$, entonces $Z_n$ converge en distribución a $N(0,1)$.
\end{block}

Sea $Y_i = \dfrac{X_i-\mu}{\sigma}$ (estandarización de cada $X_i$ individual), con $E(Y_i)=0$, $\text{Var}(Y_i)=E(Y_i^2)=1$. Entonces
\[
Z_n = \frac{1}{\sqrt n}\sum_{i=1}^n Y_i.
\]
Por independencia y la propiedad de la FGM de una suma (Clase 16, $M_{aX}(t) = M_X(at)$):
\[
M_{Z_n}(t) = \left[M_Y\!\left(\frac{t}{\sqrt n}\right)\right]^n.
\]
\end{frame}

\begin{frame}{Idea de la demostración (continuación)}
\begin{block}{Desarrollo de Taylor de $M_Y$ cerca de $0$}
Como $M_Y(0)=1$, $M_Y'(0)=E(Y)=0$, $M_Y''(0)=E(Y^2)=1$, el desarrollo de Taylor de segundo orden alrededor de $0$ da
\[
M_Y(s) \approx 1 + 0\cdot s + \frac{1}{2}s^2 + o(s^2), \quad \text{cuando } s\to 0.
\]
Sustituyendo $s = t/\sqrt n$ (que $\to 0$ cuando $n\to\infty$, para $t$ fijo):
\[
M_Y\!\left(\frac{t}{\sqrt n}\right) \approx 1 + \frac{t^2}{2n} + o\!\left(\frac1n\right).
\]
\end{block}

Elevando a la $n$ y usando el límite notable $\left(1+\dfrac{a}{n}+o(1/n)\right)^n \to e^{a}$:
\[
M_{Z_n}(t) = \left[1 + \frac{t^2}{2n} + o\!\left(\frac1n\right)\right]^n \xrightarrow[n\to\infty]{} e^{t^2/2} = M_Z(t).
\]
Por la unicidad de la FGM, $Z_n$ converge en distribución a $Z\sim N(0,1)$. (Una demostración completamente rigurosa requiere justificar el paso al límite con más cuidado —usando funciones características en vez de FGM, para garantizar existencia—, pero la idea central es exactamente esta.) $\blacksquare$
\end{frame}

\begin{frame}{Forma práctica del TCL}
\begin{alertblock}{Aproximación para $n$ grande}
Para $n$ suficientemente grande (en la práctica, suele usarse $n\geq 30$ como regla general, aunque depende de cuán ``lejos de la normal'' sea la distribución original):
\[
\overline{X}_n \ \dot\sim\ N\left(\mu, \frac{\sigma^2}{n}\right), \qquad S_n = \sum_{i=1}^n X_i \ \dot\sim\ N(n\mu, n\sigma^2),
\]
donde $\dot\sim$ significa ``se distribuye aproximadamente como''. Equivalentemente,
\[
P(\overline X_n \leq a) \approx \Phi\left(\frac{a-\mu}{\sigma/\sqrt n}\right).
\]
\end{alertblock}
\end{frame}

\section{Aplicaciones: sumas}

\begin{frame}{Ejemplo: suma de tiempos de servicio}
\begin{example}
El tiempo de atención en una caja tiene media $\mu=4$ minutos y desviación estándar $\sigma=2$ minutos (distribución no especificada). Se atienden $n=64$ clientes independientes. Aproximar la probabilidad de que el tiempo total de atención supere las $4.5$ horas ($270$ minutos).
\end{example}

\medskip
$S_{64} \dot\sim N(64\cdot 4,\ 64\cdot 4) = N(256, 256)$, con $\sqrt{256}=16$.
\[
P(S_{64} > 270) \approx P\left(Z > \frac{270-256}{16}\right) = P(Z > 0.875) \approx 1-\Phi(0.875) \approx 1-0.8092 = 0.1908.
\]
\textbf{Hay aproximadamente $19\%$ de probabilidad} de superar las $4.5$ horas.
\end{frame}

\section{Aproximación normal a la binomial}

\begin{frame}{Motivación y enunciado}
\begin{block}{Idea}
$X \sim \text{Bin}(n,p)$ puede escribirse como $X = \sum_{i=1}^n X_i$ con $X_i \sim \text{Bernoulli}(p)$ i.i.d. Por el TCL, para $n$ grande,
\[
X \ \dot\sim\ N\big(np,\; np(1-p)\big).
\]
\end{block}

\begin{alertblock}{Regla práctica de aplicabilidad}
La aproximación es razonable cuando $np \geq 5$ y $n(1-p)\geq 5$ (algunos textos piden $\geq 10$): la distribución binomial debe ser suficientemente simétrica.
\end{alertblock}

\begin{block}{Corrección por continuidad}
Como la binomial es discreta y la normal es continua, se mejora la aproximación con la \textbf{corrección por continuidad}: para $k$ entero,
\[
P(X = k) \approx P(k-0.5 < Y < k+0.5), \quad P(X\leq k)\approx P(Y \leq k+0.5), \quad Y\sim N(np,np(1-p)).
\]
\end{block}
\end{frame}

\begin{frame}{Ejemplo resuelto: aproximación normal a la binomial}
\begin{example}
Una empresa afirma que el $30\%$ de sus clientes usa la app móvil. Se encuesta a $n=200$ clientes independientes. Sea $X$ = número de clientes que usan la app. Aproximar $P(50 \leq X \leq 70)$.
\end{example}

\medskip
$X\sim\text{Bin}(200, 0.3)$: $np=60\geq 5$, $n(1-p)=140\geq 5$. $\mu = 60$, $\sigma^2 = 200(0.3)(0.7)=42$, $\sigma=\sqrt{42}\approx 6.481$.

Con corrección por continuidad:
\[
P(50\leq X\leq 70) \approx P(49.5 \leq Y \leq 70.5), \quad Y\sim N(60,42).
\]
\[
= \Phi\!\left(\frac{70.5-60}{6.481}\right) - \Phi\!\left(\frac{49.5-60}{6.481}\right) = \Phi(1.620) - \Phi(-1.620)
\]
\[
\approx 0.9474 - 0.0526 = 0.8948.
\]
\textbf{La probabilidad aproximada es $89.5\%$.}
\end{frame}

\section{Aproximación normal a la Poisson}

\begin{frame}{Motivación y enunciado}
\begin{block}{Idea}
Si $X\sim\text{Poisson}(\lambda)$ con $\lambda$ grande, $X$ puede pensarse (informalmente) como una suma de muchas variables Poisson pequeñas independientes (la Poisson es ``cerrada bajo sumas'': si $X_1\sim\text{Poisson}(\lambda_1)$, $X_2\sim\text{Poisson}(\lambda_2)$ independientes, $X_1+X_2\sim\text{Poisson}(\lambda_1+\lambda_2)$). Como $E(X)=\text{Var}(X)=\lambda$, el TCL sugiere
\[
X \ \dot\sim\ N(\lambda, \lambda), \qquad \text{razonable para } \lambda \gtrsim 10.
\]
\end{block}

\begin{alertblock}{Corrección por continuidad}
Igual que en la binomial: $P(X\leq k) \approx \Phi\!\left(\dfrac{k+0.5-\lambda}{\sqrt\lambda}\right)$.
\end{alertblock}
\end{frame}

\begin{frame}{Ejemplo resuelto: aproximación normal a la Poisson}
\begin{example}
El número de llamadas que recibe un call center en una hora sigue una distribución Poisson con $\lambda = 50$. Aproximar $P(X > 60)$.
\end{example}

\medskip
$\mu=\lambda=50$, $\sigma=\sqrt{50}\approx 7.071$. Con corrección por continuidad, $P(X>60) = P(X\geq 61) \approx P(Y > 60.5)$:
\[
P(X>60) \approx 1-\Phi\!\left(\frac{60.5-50}{7.071}\right) = 1-\Phi(1.485) \approx 1-0.9312 = 0.0688.
\]
\textbf{Hay aproximadamente $6.9\%$ de probabilidad} de recibir más de $60$ llamadas en una hora.

\medskip
\begin{block}{Comparación (referencia)}
El valor exacto (calculado con la fórmula de Poisson, sumando $P(X=61)+P(X=62)+\cdots$) es aproximadamente $0.0655$: la aproximación normal con corrección por continuidad es bastante precisa incluso con $\lambda$ moderado.
\end{block}
\end{frame}

\section{Visualización y alcance del teorema}

\begin{frame}{Visualización: la binomial se aproxima a la normal}
\begin{block}{Idea gráfica}
A medida que $n$ crece, el histograma de la distribución binomial (barras) se asemeja cada vez más a la curva de densidad normal continua. El siguiente gráfico esquematiza esta aproximación para $\text{Bin}(n,0.5)$ estandarizada.
\end{block}

\begin{center}
\begin{tikzpicture}[scale=0.85]
\begin{axis}[
  width=10cm, height=5cm,
  xlabel={valor estandarizado}, ylabel={densidad / probabilidad},
  xmin=-4, xmax=4, ymin=0, ymax=0.45,
  legend pos=north east, legend style={font=\tiny},
  ticks=none
]
\addplot[domain=-4:4, samples=100, thick, blue] {1/sqrt(2*pi)*exp(-x^2/2)};
\addplot[ycomb, thick, red, mark=*, mark size=1.2pt] coordinates {
(-3,0.011) (-2.5,0.020) (-2,0.054) (-1.5,0.099) (-1,0.157) (-0.5,0.198)
(0,0.219) (0.5,0.198) (1,0.157) (1.5,0.099) (2,0.054) (2.5,0.020) (3,0.011)
};
\legend{$N(0,1)$, $\text{Bin}(n,0.5)$ estand.}
\end{axis}
\end{tikzpicture}
\end{center}
\end{frame}

\begin{frame}{Ejemplo adicional: aproximación normal a la binomial}
\begin{example}
En una línea de producción, la probabilidad de que un artículo pase el control de calidad es $p=0.85$. Se inspeccionan $n=120$ artículos independientes. Aproximar $P(X < 95)$, donde $X$ es el número de artículos que pasan el control.
\end{example}

\medskip
$np=120(0.85)=102 \geq 5$, $n(1-p)=120(0.15)=18\geq5$: se puede aproximar. $\mu=102$, $\sigma^2=120(0.85)(0.15)=15.3$, $\sigma\approx3.912$.

Con corrección por continuidad, $P(X<95) = P(X\leq 94) \approx P(Y\leq 94.5)$:
\[
P(X<95) \approx \Phi\left(\frac{94.5-102}{3.912}\right) = \Phi(-1.917) \approx 0.0276.
\]
\textbf{Hay aproximadamente $2.76\%$ de probabilidad} de que menos de $95$ artículos pasen el control.
\end{frame}

\begin{frame}{Alcance del TCL: por qué es tan importante}
\begin{alertblock}{Universalidad}
El TCL explica por qué la distribución normal aparece con tanta frecuencia en la naturaleza y en las ciencias sociales: cualquier magnitud que resulte de \textbf{sumar o promediar} muchos efectos pequeños e independientes (errores de medición, alturas influidas por muchos genes y factores ambientales, promedios de calificaciones, etc.) tiende a distribuirse aproximadamente normal, sin importar la distribución de cada efecto individual.
\end{alertblock}

\begin{block}{Conexión con unidades futuras}
El TCL es la base teórica de la \textbf{inferencia estadística} (estimación de parámetros, intervalos de confianza, pruebas de hipótesis), que se estudiará en unidades posteriores de la materia: permite aproximar la distribución de estadísticos muestrales (como $\overline X_n$ o proporciones muestrales) por una normal, aun cuando la población de origen no lo sea.
\end{block}
\end{frame}

\section{Resumen}

\begin{frame}{Resumen}
\begin{block}{Conceptos clave}
\begin{itemize}
\item El \textbf{TCL} afirma que $Z_n = \dfrac{\overline X_n - \mu}{\sigma/\sqrt n}$ converge en distribución a $N(0,1)$, sin importar la distribución original de los $X_i$ (con media y varianza finitas).
\item Se demuestra (idea) usando que $M_{Z_n}(t) \to e^{t^2/2}$, la FGM de la normal estándar, y la propiedad de unicidad de la FGM.
\item Forma práctica: $\overline X_n \dot\sim N(\mu,\sigma^2/n)$, $S_n \dot\sim N(n\mu, n\sigma^2)$, para $n$ moderadamente grande.
\item \textbf{Aproximación normal a la binomial}: $\text{Bin}(n,p) \dot\sim N(np, np(1-p))$ si $np\geq5$ y $n(1-p)\geq5$.
\item \textbf{Aproximación normal a la Poisson}: $\text{Poisson}(\lambda)\dot\sim N(\lambda,\lambda)$ para $\lambda$ grande.
\item En ambos casos, usar \textbf{corrección por continuidad} al aproximar variables discretas con la normal continua.
\end{itemize}
\end{block}

\begin{alertblock}{Cierre de la Unidad 4}
Con esperanza, varianza, covarianza, FGM, ley de los grandes números y TCL completamos las herramientas fundamentales para el análisis probabilístico y la posterior inferencia estadística.
\end{alertblock}
\end{frame}

\end{document}
